\chapter{Introdução}
\label{cap:introducao}

\section{Contextualização}
\label{sec:contextualizacao}

Modelos de Linguagem de Grande Escala (\textit{Large Language Models}, LLMs) mudaram de forma decisiva o campo do processamento de linguagem natural (\textit{Natural Language Processing}, NLP). Antes da popularização dos LLMs, grande parte das abordagens em NLP baseava-se em \textit{pipelines} rígidos, apoiados em técnicas como \textit{tokenization}, \textit{stemming}, \textit{lemmatization}, TF-IDF e remoção de \textit{stopwords}, que transformavam texto bruto em representações numéricas esparsas, como \textit{Bag-of-Words}, para então alimentar modelos estatísticos ou algoritmos clássicos de aprendizado de máquina. Essas abordagens dependiam fortemente de engenharia manual de atributos (\textit{feature engineering}) e de ajuste fino (\textit{fine-tuning}) específico para cada conjunto de dados, o que dificultava a generalização para novos domínios.

A introdução de representações distribuídas, como o Word2Vec em 2013 \cite{mikolov2013word2vec}, começou a mudar o panorama ao permitir que modelos neurais aprendessem, de forma automática, vetores densos que capturavam as relações semânticas e sintáticas entre palavras, reduzindo a necessidade das etapas de normalização descritas anteriormente. No entanto, foi a adoção da arquitetura \textit{Transformer}, proposta por Vaswani et al. em 2017 \cite{vaswani2017attention}, que efetivamente reorganizou o campo. Modelos baseados nessa arquitetura, como o BERT \cite{devlin2019bert} e o GPT \cite{radford2018gpt}, introduziram o paradigma de pré-treinamento com grandes volumes de dados de texto não estruturados, seguido de ajuste fino em tarefas específicas com conjuntos menores e rotulados. Esse paradigma de aprendizado por transferência (\textit{transfer learning}) \cite{ruder2019transfer} permitiu que representações contextualizadas fossem aprendidas diretamente de grandes corpora não rotulados, elevando o desempenho em diversas tarefas.

Com a disponibilização comercial do ChatGPT em novembro de 2022, esses modelos deixaram de ser apenas objetos de pesquisa para se tornarem componentes centrais de sistemas que realizam tarefas de análise de sentimento, classificação de textos, extração de informações específicas de documentos e apoio à programação \cite{devlin2019bert,brown2020fewshot}. Modelos mais recentes, das famílias GPT, Claude, Gemini e Llama, são hoje integrados a assistentes virtuais, fluxos de análise e produção textual e ambientes de desenvolvimento de \textit{software}, tornando-se parte de uma infraestrutura de base em aplicações modernas nas áreas financeira, da saúde, da programação e da educação \cite{huynh2025codegen,scarlatos2025tutors}. Apesar desse avanço, todos esses modelos ainda esbarram em uma limitação estrutural fundamental, a \textbf{janela de contexto}.

A janela de contexto corresponde ao número máximo de \textit{tokens} que um modelo consegue processar como \textit{input} em uma requisição, funcionando como uma memória de trabalho disponível durante a interação. À medida que esse limite aumenta, cresce o custo computacional, pois o mecanismo de atenção apresenta complexidade quadrática $O(n^2)$ no comprimento da sequência. Isso ocorre porque cada \textit{token} é comparado com todos os demais da sequência, de modo que tanto o consumo de memória em GPUs quanto o tempo de processamento aumentam de forma acentuada conforme a janela de contexto se expande \cite{vaswani2017attention,huang2024longcontext}. Na prática, essa restrição impacta diretamente a viabilidade de aplicações que exigem interações longas com o usuário, como a análise de grandes documentos, a inspeção de grandes volumes de \textit{logs} ou o processamento de múltiplos arquivos em uma única chamada.

As janelas de contexto dos modelos baseados em \textit{Transformer} cresceram de forma expressiva ao longo do tempo. Entre 2018 e 2019, o BERT apresentava janela de contexto de 512 \textit{tokens} \cite{devlin2019bert}. Em 2023, o GPT-4 Turbo ampliou esse limite para 128.000 \textit{tokens}, e em 2024 arquiteturas como o Gemini 1.5 romperam a barreira de um milhão de \textit{tokens} \cite{reid2024gemini15}.
% [VERIFICAR REFERENCIA] A janela de 128.000 tokens do GPT-4 Turbo e a janela do Llama 3.1 nao possuem citacao. Incluir as fichas tecnicas oficiais ou remover os numeros nao verificaveis.
Nos lançamentos posteriores, o crescimento da janela parece ter perdido centralidade em favor da qualidade e da eficiência das respostas, observação que decorre do acompanhamento dos anúncios das famílias comerciais e não de uma medição sistemática.
Essa expansão, contudo, não eliminou os problemas associados ao contexto longo. Além da complexidade computacional já mencionada, fenômenos como a degradação de desempenho com o aumento do contexto (\textit{context rot}) \cite{chroma2024curse} e o colapso da atenção (\textit{rank-collapse}), em que as pontuações de atenção tendem à uniformidade, comprometem a eficácia dos modelos em cenários de contexto genuinamente extenso \cite{chen2025criticalattention}.

\subsection{A Relação Crítica entre Janela de Contexto e Desempenho}
\label{subsec:problema-critico}

Apesar da capacidade teórica de processar milhões de \textit{tokens}, os modelos modernos podem apresentar alucinações associadas ao uso ineficaz do contexto disponível. Estudos recentes \cite{chroma2024curse,chen2025criticalattention} relatam a degradação do desempenho à medida que mais informação é inserida na janela de contexto. Uma interpretação possível é que a atenção constitui um recurso finito, cujos retornos diminuem conforme novos \textit{tokens} são adicionados. Nessa leitura, cada \textit{token} acrescentado consome parte de um ``orçamento de atenção'', o que dificulta a priorização das informações mais importantes.

Além disso, surgem desafios práticos e econômicos associados ao uso de contextos extensos, descritos a seguir.

\begin{itemize}[leftmargin=1.6em]
\item \textbf{Viés \textit{lost in the middle}}. \textit{Tokens} posicionados no início ou no final da sequência tendem a receber mais atenção, enquanto aqueles situados na região central tendem a ser ignorados. Liu et al. \cite{liu2024lostmiddle} mostram que o desempenho dos LLMs segue uma curva em formato de ``U'', em que a informação no meio de grandes documentos é frequentemente desprezada, o que pode ocasionar alucinações.

\item \textbf{Custo computacional e econômico}. Como o processamento de \textit{tokens} é o principal componente de custo em servidores, o uso de contextos muito longos pode tornar determinadas aplicações economicamente inviáveis, exigindo técnicas de otimização como a compressão de \textit{prompt} \cite{jiang2023llmlingua}.

\item \textbf{Latência}. Devido à complexidade $O(n^2)$ da atenção \cite{vaswani2017attention}, contextos extensos aumentam drasticamente o tempo de \textit{prefill} e de geração, elevando a latência em aplicações interativas \cite{aman2024latency}.
\end{itemize}

\section{Problema de Pesquisa}
\label{sec:problema-pesquisa}

Os avanços nas arquiteturas e nas técnicas de extensão de contexto ainda não superam a limitação prática do uso de LLMs em cenários cotidianos. Os principais modelos, acessados por interface ou por API, como ChatGPT e Claude, impõem limites ao comprimento da entrada e continuam sujeitos a alucinações, especialmente em contextos extensos. Para contornar essa limitação, tornou-se comum a adoção de mecanismos de memória externa e de geração aumentada por recuperação (\textit{Retrieval-Augmented Generation}, RAG). Contudo, na prática, mecanismos de busca vetorial aproximam similaridade por distância em um espaço de \textit{embeddings}, o que não garante, em todos os casos, que os trechos recuperados sejam os mais relevantes para a tarefa. Yu et al. \cite{yu2024defenserag} apresentam uma discussão mais ampla sobre os limites das aplicações de RAG.

Um ponto pouco explorado é que a maioria dessas técnicas trata a compressão de forma \textbf{uniforme}, aplicando um único método, e frequentemente uma única taxa de compressão, a todo o \textit{prompt}. Essa escolha desconsidera a estrutura interna do \textit{prompt}. Instruções, perguntas, restrições de formato, evidências recuperadas, exemplos e trechos de fundo não carregam o mesmo valor para a tarefa, nem toleram o mesmo grau de compressão. Uma taxa uniforme pode comprimir demais uma região crítica e comprimir de menos uma região redundante.

Diante desse cenário, este trabalho parte da seguinte pergunta central.

\begin{quote}
\itshape Como estender, de forma não invasiva, a quantidade de informação útil que um LLM acessado via API consegue aproveitar dentro de sua janela de contexto, selecionando e comprimindo seletivamente as regiões do \textit{prompt} de modo a preservar a informação mais relevante para a tarefa sob uma restrição de orçamento de \textit{tokens}?
\end{quote}

Diferentemente de abordagens que modificam diretamente o mecanismo de atenção ou os \textit{position embeddings}, este estudo propõe um \textit{framework} que atua sobre a entrada do modelo. A ideia central é tratar a compressão de \textit{prompt} não como uma redução uniforme de comprimento, mas como um problema de \textbf{alocação de orçamento}, distribuindo a pressão de compressão entre as regiões do \textit{prompt} segundo seu valor informacional e seu risco de perda, e formulando a seleção de compressores como um problema da mochila de múltipla escolha.

\section{Objetivos}
\label{sec:objetivos}

\subsection{Objetivo Global}
\label{subsec:objetivo-global}

Propor um \textit{framework} não invasivo para extensão da janela de contexto em LLMs que formule a compressão de \textit{prompt} como um problema de alocação de orçamento, mais especificamente um problema da mochila de múltipla escolha. Para cada região do \textit{prompt}, o \textit{framework} seleciona a estratégia de compressão que maximiza a informação preservada sob uma restrição de orçamento de \textit{tokens}, sem retreinamento nem modificação arquitetural do modelo base.

\subsection{Objetivos Específicos}
\label{subsec:objetivos-especificos}

\begin{enumerate}[leftmargin=1.8em]
\item Revisar sistematicamente as estratégias contemporâneas de extensão de contexto que operam \textbf{sem} os três elementos a seguir.
    \begin{enumerate}[label=(\alph*)]
\item retreinamento do modelo;
\item modificação dos mecanismos de atenção;
\item alteração dos sistemas de \textit{Position Embeddings};
    \end{enumerate}

\item Avaliar empiricamente as estratégias não invasivas mais promissoras em \textit{benchmarks} especializados em contexto longo, medindo acurácia e latência em modelos de diferentes escalas e arquiteturas;

\item Formalizar a seleção de compressores por região do \textit{prompt} como um problema da mochila de múltipla escolha, com restrições de orçamento e de compatibilidade entre técnicas;

\item Especificar um conjunto de compressores compatíveis (\textit{text-in}/\textit{text-out}) e uma função de importância por partição que combine relevância à tarefa, densidade de informação e saliência posicional;

\item Discutir métricas de avaliação que considerem não apenas o desempenho final, mas também a preservação de informação, a fidelidade semântica e o risco de alucinação por perda de conteúdo.
\end{enumerate}

\section{Justificativa e Contribuições}
\label{sec:justificativa}

A pesquisa sobre extensão de contexto em LLMs tem se concentrado predominantemente em intervenções arquiteturais diretas. O mecanismo de auto-atenção, proposto por Vaswani et al. \cite{vaswani2017attention} e consolidado em modelos de escala cada vez maior, como o GPT-3 \cite{brown2020fewshot}, tornou-se o elemento central para a captura de dependências de longo alcance, inspirando uma linha de pesquisa dedicada a modificações estruturais do \textit{Transformer}. Diversos \textit{surveys} consolidam essa ênfase arquitetural, categorizando técnicas que alongam o contexto por meio de ajustes no \textit{Transformer}, otimizações do mecanismo de atenção e esquemas especializados de codificação posicional \cite{huang2024longcontext,wang2024beyondlimits,zhao2023survey,naveed2023overview,wan2023efficient,zhuang2023efficienttraining}. Embora essas abordagens demonstrem ganhos teóricos significativos, sua aplicação prática frequentemente esbarra em limitações de acesso, pois a maioria dos usuários de LLMs interage com modelos via API ou interfaces de chat, sem controle sobre a arquitetura interna ou capacidade de retreinamento.

Em contraste, este trabalho adota estratégias \textbf{não invasivas}, isto é, técnicas aplicáveis sobre modelos de API sem necessidade de retreinamento, \textit{fine-tuning} ou modificação de componentes arquiteturais. Essas estratégias incluem compressão de \textit{prompts} \cite{jiang2023llmlingua,gilbert2023semanticcompression}, segmentação e processamento sequencial de janelas de contexto \cite{kamradt2023workaround}, arquiteturas de memória externa persistente \cite{packer2024memgpt,xiao2024infllm} e recuperação aumentada de informações externas \cite{yu2024defenserag,lewis2020rag,sharma2025ragsurvey}. Trabalhos recentes têm refinado essas técnicas com abordagens como o ganho de informação relevante para seleção de trechos \cite{pickett2024rig}, a integração de grafos de conhecimento esparsos \cite{alonso2024mixturepageranks} e a otimização de memória em tempo de execução \cite{kwon2023pagedattention}. Choi et al. \cite{choi2025loraragg} mostram, por exemplo, que a combinação de RAG com adaptação de baixo \textit{rank} (LoRA) pode melhorar a precisão de recuperação ao mesmo tempo em que reduz custos computacionais, o que sugere o potencial de técnicas não invasivas em aplicações especializadas.

Além disso, este trabalho propõe um diferencial metodológico. Em vez de tratar a compressão de \textit{prompt} como redução uniforme de \textit{tokens}, introduz-se a noção de \textbf{densificação informacional orientada à tarefa}, na qual a compressão é entendida como a transformação de um contexto longo em uma representação menor que preserva a informação utilizável para a tarefa, sob restrições operacionais e epistêmicas. A partir dessa noção, a seleção do compressor de cada região do \textit{prompt} é formulada como um problema da mochila de múltipla escolha, permitindo que o \textit{framework} funcione como uma camada de decisão acima dos compressores já existentes.

\subsection{Principais Contribuições}
\label{subsec:contribuicoes}

Este trabalho oferece as seguintes contribuições.

\begin{enumerate}[leftmargin=1.8em]
\item Uma revisão sistemática de técnicas \textbf{não invasivas} de extensão de contexto, categorizando abordagens aplicáveis sem acesso aos parâmetros ou à arquitetura interna do modelo, organizada em uma taxonomia por lógica operacional;

\item Uma avaliação empírica comparativa conduzida em duas fases, uma preliminar sobre modelos servidos por APIs externas e outra local, que compara nove técnicas de tratamento de contexto em cinco modelos de arquiteturas distintas e dois orçamentos, reportando acurácia, latência e taxa de compressão efetiva em \textit{benchmarks} de contexto longo, com a fidelidade semântica medida na etapa de ampliação dedicada aos quatro compressores;

\item Um \textit{framework} de \textbf{compressão seletiva via problema da mochila}, que atua como uma camada de decisão acima de compressores existentes, alocando a pressão de compressão entre as regiões do \textit{prompt} de acordo com seu valor informacional e seu risco de perda;

\item Uma avaliação do \textit{framework} em dois orçamentos de contexto, que caracteriza o regime em que a alocação seletiva supera a compressão uniforme, acompanhada de um plano experimental reprodutível para as demais configurações, com \textit{benchmarks}, \textit{baselines} e estudos de ablação orientados por métricas de preservação de informação e não apenas de desempenho final.
\end{enumerate}

\section{Estrutura da Dissertação}
\label{sec:estrutura}

Esta dissertação está organizada em cinco capítulos, estruturados de forma a conduzir o leitor desde os fundamentos técnicos até a formulação e o plano de avaliação do \textit{framework} proposto.

O \textbf{Capítulo~\ref{cap:introducao}} contextualiza o problema da janela de contexto em LLMs, apresenta os desafios práticos de custo, latência e degradação de desempenho em contextos extensos, formula a pergunta de pesquisa central, define os objetivos e apresenta as contribuições esperadas.

O \textbf{Capítulo~\ref{cap:fundamentos}} revisa os princípios técnicos dos LLMs, que abrangem a arquitetura \textit{Transformer}, o mecanismo de atenção, a complexidade computacional e o funcionamento da janela de contexto. Em seguida, discute as limitações práticas dos modelos em contextos longos, incluindo fenômenos como \textit{context rot}, \textit{lost in the middle} e colapso de atenção, e propõe uma taxonomia das estratégias não invasivas de extensão de contexto.

O \textbf{Capítulo~\ref{cap:metodologia}} apresenta a validação empírica das técnicas não invasivas revisadas, com seleção da abordagem de melhor desempenho segundo \textit{benchmarks} especializados, e formaliza a arquitetura do \textit{framework} proposto, descrevendo a compressão seletiva de contexto como um problema da mochila de múltipla escolha, com sua função de importância, seu modelo de qualidade, suas restrições de compatibilidade e seu algoritmo de solução. O capítulo encerra com as perguntas de pesquisa, as hipóteses de trabalho e o plano experimental que orientam a avaliação.

O \textbf{Capítulo~\ref{cap:resultados}} apresenta e discute os resultados obtidos, analisa a relação entre qualidade, compressão e latência entre modelos e \textit{benchmarks}, caracteriza as opções de compressão selecionadas para o \textit{framework} e reporta a avaliação do resolvedor sob dois orçamentos de contexto, delimitando o alcance das evidências obtidas.

O \textbf{Capítulo~\ref{cap:conclusao}} sintetiza os principais achados, retoma as contribuições e discute as implicações práticas e teóricas do \textit{framework}. Identifica limitações metodológicas, desafios de escalabilidade e direções para investigações futuras, incluindo a execução completa do plano experimental, o alinhamento da função de valor ao desempenho final, a integração com modelos multimodais e a otimização de latência.
